This paper formulates consciousness not as an internally generated object, but as a state that opens temporarily between processes operating at different time scales. The central claim is that consciousness, in this descriptive layer, is not treated as a specific brain region, representation, or output, but as a log-referenceable relation established between fast predictive processes and slow sedimentary memory processes.

To describe this relation, we introduce internal delay or buffering as $\tau$, an effective torque-like quantity as $\mathcal{T}$, and information potential as $\Phi$. Speed difference is written $\Delta v = v_f - v_s$ and coupling viscosity as $\mu$. Libet-type experiments, split-brain experiments, qualia, measurement, and non-localization are all reinterpreted from the perspective of inter-scale transformation and incomplete closure.

The paper further redescribes log referencing as non-closed information flow along a meaning gradient. The IFGT variables $I$, $J$, and $\Phi$ are used here in a projected, cognitive-window sense, not as a direct identification of consciousness with a general information field.
